% GENERATED FILE. Edit canonical lessons, then run npm run generate:books.
% canonical-source-hash: fnv1a64:0f48c34343cc59f6
% canonical-lessons: FR-C215-grippe, FR-C215-ordonnance, FR-C215-comprime, FR-C215-sirop, FR-C215-pommade

\chapter{Flu, Prescription, Tablet, Syrup, Ointment}
\label{ch:fr-flu-prescription-tablet-syrup-ointment}
\hlchaptermodality{215}

\begin{chapteropening}
\textbf{By the end of this chapter:} \emph{I can say flu, a prescription, a tablet, a syrup and an ointment.}

Complete the last lesson of chapter 215: I can say flu, a prescription, a tablet, a syrup and an ointment.
\end{chapteropening}

\section[la grippe]{la grippe — flu}

\label{lesson:FR-C215-grippe}



\begin{quote}
Before the new one: say the French for to express, then the French for to put on weight.
\end{quote}

\subsection*{You'll want to know: la grippe}
\textbf{la grippe} — \textquotedblleft{}flu\textquotedblright{}.

\begin{grammarlens}[title={What you've built}]
One more word for this chapter.
\end{grammarlens}

\subsection*{Your turn}
\begin{itemize}
\raggedright
  \item \emph{Say it:} \emph{la grippe}
  \item \emph{Say it:} \emph{la grippe}, once more
  \item \emph{Say it:} say \emph{la grippe}
  \item \emph{Recall:} say the French for to express, then the French for to put on weight, then say \emph{la grippe} again
\end{itemize}

\begin{tcolorbox}[breakable,colback=teal!4,colframe=teal!35!black,title={Before you move on}]
What does la grippe mean? (\textquotedblleft{}Flu\textquotedblright{}.) Say it once more.
\end{tcolorbox}
\section[l'ordonnance]{l'ordonnance — a prescription}

\label{lesson:FR-C215-ordonnance}



\begin{quote}
Before the new one: say the French for to put on weight, then the French for flu.
\end{quote}

\subsection*{You'll want to know: l'ordonnance}
\textbf{l'ordonnance} — \textquotedblleft{}a prescription\textquotedblright{}.

\begin{grammarlens}[title={What you've built}]
One more word for this chapter.
\end{grammarlens}

\subsection*{Your turn}
\begin{itemize}
\raggedright
  \item \emph{Say it:} \emph{l'ordonnance}
  \item \emph{Say it:} \emph{l'ordonnance}, once more
  \item \emph{Say it:} say \emph{l'ordonnance}
  \item \emph{Recall:} say the French for to put on weight, then the French for flu, then say \emph{l'ordonnance} again
\end{itemize}

\begin{tcolorbox}[breakable,colback=teal!4,colframe=teal!35!black,title={Before you move on}]
What does l'ordonnance mean? (\textquotedblleft{}A prescription\textquotedblright{}.) Say it once more.
\end{tcolorbox}
\section[le comprimé]{le comprimé — a tablet (medicine)}

\label{lesson:FR-C215-comprime}



\begin{quote}
Before the new one: say the French for flu, then the French for a prescription.
\end{quote}

\subsection*{You'll want to know: le comprimé}
\textbf{le comprimé} — \textquotedblleft{}a tablet (medicine)\textquotedblright{}.

\begin{grammarlens}[title={What you've built}]
One more word for this chapter.
\end{grammarlens}

\subsection*{Your turn}
\begin{itemize}
\raggedright
  \item \emph{Say it:} \emph{le comprimé}
  \item \emph{Say it:} \emph{le comprimé}, once more
  \item \emph{Say it:} say \emph{le comprimé}
  \item \emph{Recall:} say the French for flu, then the French for a prescription, then say \emph{le comprimé} again
\end{itemize}

\begin{tcolorbox}[breakable,colback=teal!4,colframe=teal!35!black,title={Before you move on}]
What does le comprimé mean? (\textquotedblleft{}A tablet (medicine)\textquotedblright{}.) Say it once more.
\end{tcolorbox}
\section[le sirop]{le sirop — a syrup}

\label{lesson:FR-C215-sirop}



\begin{quote}
Before the new one: say the French for a prescription, then the French for a tablet.
\end{quote}

\subsection*{You'll want to know: le sirop}
\textbf{le sirop} — \textquotedblleft{}a syrup\textquotedblright{}.

\begin{grammarlens}[title={What you've built}]
One more word for this chapter.
\end{grammarlens}

\subsection*{Your turn}
\begin{itemize}
\raggedright
  \item \emph{Say it:} \emph{le sirop}
  \item \emph{Say it:} \emph{le sirop}, once more
  \item \emph{Say it:} say \emph{le sirop}
  \item \emph{Recall:} say the French for a prescription, then the French for a tablet, then say \emph{le sirop} again
\end{itemize}

\begin{tcolorbox}[breakable,colback=teal!4,colframe=teal!35!black,title={Before you move on}]
What does le sirop mean? (\textquotedblleft{}A syrup\textquotedblright{}.) Say it once more.
\end{tcolorbox}
\section[la pommade]{la pommade — an ointment}

\label{lesson:FR-C215-pommade}



\begin{quote}
Before the new one: say the French for a tablet, then the French for a syrup.
\end{quote}

\subsection*{You'll want to know: la pommade}
\textbf{la pommade} — \textquotedblleft{}an ointment\textquotedblright{}.

\begin{grammarlens}[title={What you've built}]
One more word for this chapter.
\end{grammarlens}

\subsection*{Your turn}
\begin{itemize}
\raggedright
  \item \emph{Say it:} \emph{la pommade}
  \item \emph{Say it:} \emph{la pommade}, once more
  \item \emph{Say it:} say \emph{la pommade}
  \item \emph{Recall:} say the French for a tablet, then the French for a syrup, then say \emph{la pommade} again
\end{itemize}

\begin{tcolorbox}[breakable,colback=teal!4,colframe=teal!35!black,title={Before you move on}]
What does la pommade mean? (\textquotedblleft{}An ointment\textquotedblright{}.) Say it once more.
\end{tcolorbox}
